% =====================================================================
\chapter{AutoML y búsqueda de arquitecturas}
\label{tema:3}
% =====================================================================

Elegir un modelo, sus hiperparámetros y su preprocesado es también un problema de optimización sobre programas: cada configuración es una canalización ejecutable cuyo valor se mide con validación. Este tema formaliza ese problema siguiendo el libro editado por Hutter, Kotthoff y Vanschoren~\cite{automl} (cap.~1 «Hyperparameter Optimization» y cap.~3 «Neural Architecture Search», de Elsken, Metzen y Hutter~\cite{elsken}) y los artículos originales: Auto-WEKA~\cite{autoweka}, búsqueda aleatoria~\cite{bergstra}, búsqueda de arquitecturas con refuerzo~\cite{zoph}, con evolución~\cite{real}, con pesos compartidos~\cite{enas} y diferenciable~\cite{darts}, y la línea base de Li y Talwalkar~\cite{lt}. La multifidelidad sigue a Jamieson y Talwalkar~\cite{sh}, a Li \emph{et al.}~\cite{hyperband} y a la §19.4 de Zhang \emph{et al.}~\cite{d2l}. Complementa el tema~6 del módulo~25, que estudió el ajuste de hiperparámetros con optimización bayesiana y Hyperband; aquí el énfasis está en el espacio de búsqueda y en el coste.

\section{AutoML como optimización sobre programas}

\begin{definicion}[Selección conjunta de algoritmo e hiperparámetros]\label{def:cash}
Se consideran $R$ algoritmos de aprendizaje, $A^{(1)},\dots,A^{(R)}$, cada uno con un espacio de hiperparámetros $\Lambda^{(j)}$, y sea $D$ un conjunto de datos dividido en $k$ pares $(D^{(i)}_{\mathrm{tr}},D^{(i)}_{\mathrm{va}})$, cada uno formado por una parte de entrenamiento y otra de validación. El problema CASH (\emph{combined algorithm selection and hyperparameter optimization})~\cite{autoweka} es
\[
(j^{*},\lambda^{*})\in\argmin_{j\in\{1,\dots,R\},\ \lambda\in\Lambda^{(j)}}\ \frac1k\sum_{i=1}^{k}\mathcal{L}\bigl(A^{(j)}_{\lambda},D^{(i)}_{\mathrm{tr}},D^{(i)}_{\mathrm{va}}\bigr),
\]
donde $\mathcal{L}$ es la pérdida de validación del modelo entrenado con $D^{(i)}_{\mathrm{tr}}$.
\end{definicion}

El espacio de la \cref{def:cash} es \emph{condicional}: la profundidad de un árbol solo existe si se elige un árbol. Se representa como un único espacio jerárquico cuya raíz es el hiperparámetro «algoritmo», de modo que métodos como la optimización bayesiana con bosques aleatorios (SMAC) o con estimadores de densidad en árbol (TPE) lo tratan sin cambios. Si además se eligen el preprocesado, la selección de variables y la combinación de modelos, la configuración es un programa completo; TPOT, por ejemplo, busca esas canalizaciones con programación genética, como en el \cref{tema:1}~\cite{automl}.

\begin{observacion}[Por qué la búsqueda aleatoria es una línea base seria]\label{obs:aleatoria}
Si la pérdida depende sobre todo de unos pocos hiperparámetros, una rejilla desperdicia evaluaciones repitiendo los mismos valores de los importantes, mientras que $n$ puntos aleatorios prueban $n$ valores distintos de cada uno~\cite{bergstra}. Además, la probabilidad de que al menos uno de $n$ puntos aleatorios caiga en la fracción $q$ mejor del espacio es $1-(1-q)^{n}$, independiente de la dimensión: con $q=0{,}05$ y $n=60$ vale $0{,}954$ (módulo~25, tema~6).
\end{observacion}

\section{Búsqueda de arquitecturas neuronales}

Elsken, Metzen y Hutter organizan la búsqueda de arquitecturas (NAS) en tres dimensiones~\cite{elsken}: el \emph{espacio de búsqueda} (qué arquitecturas pueden representarse), la \emph{estrategia de búsqueda} (cómo se explora) y la \emph{estimación del rendimiento} (cómo se evalúa una arquitectura sin entrenarla del todo).

\begin{definicion}[Espacio de celdas]\label{def:celdas}
Una \emph{celda} es un grafo dirigido acíclico con dos nodos de entrada (las salidas de las dos celdas anteriores) y $B$ nodos intermedios. El nodo intermedio $b$ elige dos entradas entre los $b+1$ nodos anteriores y aplica a cada una una operación de un conjunto $\mathcal{O}$ (convoluciones, agrupamientos, identidad\dots); su salida es la suma de ambas. La red se construye apilando copias de una celda normal y de una celda de reducción.
\end{definicion}

Si las dos elecciones de cada nodo son ordenadas y pueden repetirse, el número de celdas distintas es $\prod_{b=1}^{B}\bigl((b+1)\abs{\mathcal{O}}\bigr)^{2}$: con $B=5$ y $\abs{\mathcal{O}}=8$, unas $5{,}6\cdot10^{14}$. La celda reduce el espacio respecto a una red arbitraria, pero sigue siendo inabarcable por enumeración.

\emph{Estrategias de búsqueda.} Zoph y Le entrenaron con refuerzo una red recurrente que genera descripciones de arquitecturas y obtuvieron un $3{,}65\,\%$ de error en CIFAR-10~\cite{zoph}. Real \emph{et al.} usaron evolución por torneo con \emph{envejecimiento} (se elimina el individuo más antiguo, no el peor), y en una comparación controlada con el mismo hardware la evolución llegó antes a buenos resultados que el refuerzo~\cite{real}. Liu, Simonyan y Yang hicieron el espacio continuo:

\begin{definicion}[Relajación continua de DARTS]\label{def:darts}
En DARTS~\cite{darts}, cada arista $(i,j)$ de la celda aplica una \emph{mezcla} de todas las operaciones,
\[
\bar o^{(i,j)}(x)=\sum_{o\in\mathcal{O}}\frac{\exp\bigl(\alpha^{(i,j)}_{o}\bigr)}{\sum_{o'\in\mathcal{O}}\exp\bigl(\alpha^{(i,j)}_{o'}\bigr)}\,o(x),
\]
y los parámetros de arquitectura $\alpha$ se optimizan por descenso de gradiente en un problema binivel: $\min_{\alpha}\mathcal{L}_{\mathrm{va}}\bigl(w^{*}(\alpha),\alpha\bigr)$ con $w^{*}(\alpha)=\argmin_{w}\mathcal{L}_{\mathrm{tr}}(w,\alpha)$. Al terminar, cada arista conserva la operación de mayor $\alpha$.
\end{definicion}

\emph{Estimación del rendimiento.} Entrenar cada candidata hasta el final es lo que hace prohibitiva la búsqueda. ENAS comparte los pesos entre todas las arquitecturas del espacio (cada candidata es un subgrafo de un modelo grande) y lo abarata en unas mil veces en horas de GPU respecto a la búsqueda original~\cite{enas}. Li y Talwalkar mostraron que la búsqueda aleatoria con parada temprana iguala a ENAS en PTB y CIFAR-10, y con pesos compartidos lo supera, y señalaron problemas de reproducibilidad en la literatura~\cite{lt}. La lección práctica es la del \cref{tema:1}: compárese siempre con una búsqueda aleatoria del mismo presupuesto.

\section{Multifidelidad y coste}

\begin{definicion}[Fidelidad]\label{def:fidelidad}
Una \emph{fidelidad} es un recurso $r\in[r_{\min},r_{\max}]$ que controla el coste y la precisión de la evaluación de una configuración $\lambda$: épocas de entrenamiento, fracción de los datos, resolución de las imágenes o número de árboles. Se escribe $f(\lambda,r)$ para la pérdida de validación con presupuesto $r$; la de interés es $f(\lambda,r_{\max})$.
\end{definicion}

\begin{definicion}[Halving sucesivo]\label{def:sh}
Dados $N$ configuraciones, un factor $\eta\ge2$ y un presupuesto mínimo $r_{\min}$, el \emph{halving sucesivo}~\cite{sh} evalúa todas con $r_0=r_{\min}$, conserva la fracción $1/\eta$ mejor, multiplica el presupuesto por $\eta$ y repite: en el escalón $i$ quedan $n_i=N/\eta^{i}$ configuraciones con presupuesto $r_i=r_{\min}\eta^{i}$, hasta $r_K=r_{\max}$.
\end{definicion}

\begin{proposicion}[Coste del halving sucesivo]\label{prop:coste-sh}
Si $N$ es múltiplo de $\eta^{K}$ y el coste de una evaluación es proporcional a su presupuesto, cada escalón cuesta $N r_{\min}$ y el total es $(K+1)Nr_{\min}$, frente a $Nr_{\max}=N r_{\min}\eta^{K}$ de evaluar todas con el presupuesto completo. El cociente es $(K+1)/\eta^{K}$.
\end{proposicion}

\begin{proposicion}[Cuándo acierta el halving sucesivo]\label{prop:acierto-sh}
Sea $\lambda^{\dagger}$ la mejor configuración a fidelidad completa. Si en cada escalón $i<K$ la ordenación por $f(\cdot,r_i)$ coloca a $\lambda^{\dagger}$ entre las $n_{i+1}$ primeras de las que quedan, el halving sucesivo devuelve $\lambda^{\dagger}$. Si las curvas de aprendizaje se cruzan (una configuración peor al principio es mejor al final), puede descartarla.
\end{proposicion}

La hipótesis de la \cref{prop:acierto-sh} se puede comprobar en un piloto con la correlación de rangos de Spearman entre las pérdidas a baja y a alta fidelidad. Hyperband~\cite{hyperband} protege frente a una mala elección de $r_{\min}$ ejecutando varias rondas de halving sucesivo con distinta agresividad (módulo~25, tema~6). En la práctica, el coste no es exactamente proporcional al presupuesto: cargar los datos, compilar el modelo o arrancar una máquina son costes fijos que reducen el ahorro de la \cref{prop:coste-sh}.

\section{Caso de estudio: AutoML del clasificador de averías}

\addtocounter{definicion}{1}%

\begin{table}[h]
\centering\small
\caption{Tecnologías del tema~3.}\label{tab:tec3}
\begin{tabularx}{\textwidth}{@{}>{\raggedright\arraybackslash}p{2.5cm}>{\raggedright\arraybackslash}X>{\raggedright\arraybackslash}p{2.1cm}>{\raggedright\arraybackslash}X@{}}
\toprule
Herramienta & Para qué sirve & Licencia o tipo & Cuándo elegirla \\
\midrule
Optuna~\cite{optuna} & optimización de hiperparámetros con espacios definidos en el propio código (TPE y otros) y poda con halving sucesivo o Hyperband & abierta (MIT) & ajuste de cualquier modelo en Python con poco código \\
Ray Tune~\cite{ray} & ajuste distribuido en un clúster con planificadores de multifidelidad & abierta (Apache~2.0) & muchas pruebas en paralelo sobre varias máquinas \\
SMAC3 & configuración algorítmica con bosques aleatorios como sustituto & abierta (BSD-3) & espacios condicionales, configurar resolvedores \\
auto-sklearn, AutoGluon & AutoML completo para datos tabulares (selección, ajuste, conjuntos) & abiertas (BSD-3; Apache~2.0) & primera línea base fuerte en datos tabulares \\
NNI & búsqueda de arquitecturas y de hiperparámetros & abierta (MIT) & NAS con estrategias variadas sobre PyTorch \\
\bottomrule
\end{tabularx}
\end{table}

\addtocounter{definicion}{1}%

